% ID: SEL-D0055-OA-E01
% ESTADO: validada. S056.
\textbf{a)} El recinto se sitúa por debajo de $y=6x+9$ y $y=(13-2x)/5$, y por encima de $y=(2x-5)/3$. Las fronteras están incluidas. No se impone que las coordenadas sean positivas.
\begin{center}
\setlength{\unitlength}{0.9mm}
\begin{picture}(97,61)(-39,-29)
\qbezier(-32,0)(9,0)(50,0)
\qbezier(0,-25)(0,0)(0,25)
\put(51,0){$x$}
\put(2,25){$y$}
\qbezier(-20,-21)(-15,0)(-10,21)
\qbezier(-10,21)(15,14)(40,7)
\qbezier(40,7)(10,-7)(-20,-21)
\put(-37,-24){\scriptsize $(-2,-3)$}
\put(-25,23){\scriptsize $(-1,3)$}
\put(41,8){\scriptsize $(4,1)$}
\end{picture}
\end{center}
\textbf{b)} Igualando las dos cotas superiores obtenemos $(-1,3)$. Al igualar la cota inferior con cada una de ellas salen $(-2,-3)$ y $(4,1)$. Los tres puntos satisfacen las inecuaciones.

\textbf{c)} En $(-2,-3),(-1,3),(4,1)$, los valores de $F$ son $3,-6,13$, respectivamente. Por linealidad, el máximo es $\boxed{13}$ sólo en $(4,1)$ y el mínimo es $\boxed{-6}$ sólo en $(-1,3)$, ya que en cada caso sólo un vértice alcanza el valor extremo.
